\documentclass[11pt]{article}
\usepackage[T1]{fontenc}
\usepackage[utf8]{inputenc}
\usepackage[french]{babel}
\usepackage{newpxtext,newpxmath}
\usepackage[a4paper,top=1.9cm,bottom=2.0cm,left=2.1cm,right=2.1cm]{geometry}
\usepackage{microtype}
\usepackage{xcolor}
\usepackage{titlesec}
\usepackage[most]{tcolorbox}
\usepackage{fancyhdr}
\usepackage{listings}

\definecolor{efenavy}{HTML}{1F3A5F}
\definecolor{efeteal}{HTML}{0E7A6E}
\definecolor{efeamber}{HTML}{B25A1E}
\definecolor{efeplum}{HTML}{7A4EA8}
\definecolor{efegray}{HTML}{54606E}
\definecolor{boxgray}{HTML}{F3F5F8}
\definecolor{boxamber}{HTML}{FBF3EA}
\definecolor{boxedge}{HTML}{D9E0EA}
\definecolor{codebg}{HTML}{F5F7FA}

\lstdefinestyle{rs}{basicstyle=\ttfamily\footnotesize,backgroundcolor=\color{codebg},
  keywordstyle=\color{efenavy}\bfseries,commentstyle=\color{efegray}\itshape,
  stringstyle=\color{efeteal},breaklines=true,columns=fullflexible,keepspaces=true,
  xleftmargin=8pt,framexleftmargin=8pt,aboveskip=5pt,belowskip=3pt,
  morekeywords={pub,struct,enum,fn,impl,let,match,Result,Ok,Err,const,for,assert,use,mod,None,Some}}
\lstset{style=rs}

\titleformat{\section}{\color{efenavy}\large\bfseries}{}{0pt}{}
\titlespacing{\section}{0pt}{11pt}{4pt}

\setlength{\parskip}{3pt}\setlength{\parindent}{0pt}\setlength{\emergencystretch}{3em}

\pagestyle{fancy}\fancyhf{}\renewcommand{\headrulewidth}{0.3pt}
\lhead{\small\color{efegray}reg-evidence — outillé $\neq$ conforme}
\rhead{\small\color{efegray}2026-07-19}
\cfoot{\small\color{efegray}\thepage}

\newtcolorbox{these}{colback=boxgray,colframe=boxedge,boxrule=0.5pt,arc=2.5pt,left=8pt,right=8pt,top=5pt,bottom=5pt}
\newtcolorbox{cadre}{colback=boxamber,colframe=efeamber!55,boxrule=0.6pt,arc=2.5pt,left=8pt,right=8pt,top=5pt,bottom=5pt}

\begin{document}

\begin{center}
{\LARGE\bfseries\color{efenavy} reg-evidence — outillé $\neq$ conforme}\\[3pt]
{\large\color{efegray}\itshape Le reçu re-vérifiable, et l'artefact réglementaire qui refuse
par construction de dire « conforme »}\\[5pt]
{\color{efegray}Rodrigue Pekar \quad$\cdot$\quad Note technique \quad$\cdot$\quad 2026-07-19}
\end{center}
\vspace{2pt}

\begin{these}
Tout outil de conformité vend une \textbf{coche verte} : « votre système est conforme ».
\texttt{reg-evidence} fait l'inverse — et il l'impose par le \textbf{type} et par la
\textbf{cryptographie}, pas par un discours. L'artefact réglementaire qu'il produit n'a
\emph{structurellement aucun champ} où écrire « conforme » : seulement de la \textbf{couverture}
(couvert / partiel / hors-périmètre). Et il est dérivé d'un \textbf{reçu scellé} qu'un tiers
\emph{sceptique} re-prouve sans jamais faire confiance au producteur. La conformité reste la
conclusion du juriste ; l'outil ne fournit que de l'évidence re-vérifiable. Refuser le verdict
est plus dur à vendre qu'une coche verte — et bien plus solide.
\end{these}

\section{1 \;\textbar\; La coche verte que tout le monde vend}
Le marché de la conformité repose sur une phrase : « conforme SOC 2 », « prêt RGPD »,
« AI-Act-ready ». C'est un \emph{verdict} — or un verdict de conformité dépend de la
classification d'un système précis, de son déploiement, du conseil juridique du client. Un
outil ne peut pas le prononcer honnêtement. La plupart le prononcent quand même, en petits
caractères. \texttt{reg-evidence} prend la position inverse et la rend \emph{impossible à
contourner} : il n'y a nulle part, dans le type de sortie, où écrire « conforme ». Ce document
décrit les deux couches qui rendent ce refus non pas rhétorique mais \textbf{mécanique}.

\section{2 \;\textbar\; Couche 0 — le reçu re-vérifiable (\texttt{provenance-dag})}
Chaque décision, chaque appel, chaque hop est un \textbf{nœud} scellé dans un DAG. Le sceau
lie le nœud à toute la chaîne qui précède :
\[
\text{sceau}_i = \mathrm{HMAC}\big(K_{\text{tenant}},\ \text{chaîne}_{i-1}\ \Vert\ \mathrm{canonical}(\text{n\oe ud}_i)\big)
\]
\vspace{-6pt}
\[
\text{avec}\quad \text{chaîne}_0 = \mathrm{SHA256}(\text{GENESIS}).
\]
\texttt{append()} (\texttt{provenance-dag/src/dag.rs:83}, sceau ligne 95) scelle chaque nœud et
\emph{refuse à l'écriture} les incohérences de structure : référence causale pendante
(test \texttt{dag.rs:343}) et identifiant dupliqué (test \texttt{dag.rs:361}). Rien de mal formé
n'entre dans la chaîne.

\textbf{La re-preuve sans confiance.} Un tiers qui possède seulement les \emph{records} et la
clé de tenant appelle \texttt{reprove\_dag} (\texttt{dag.rs:167}) : il \textbf{re-dérive toute la
chaîne} depuis GENESIS et recompte chaque sceau, comparé en \emph{temps constant}
(\texttt{ct\_eq}, \texttt{dag.rs:188}). Il ne fait confiance à \emph{rien} de ce que le producteur
affirme — il recalcule. \texttt{verify\_against\_root} (\texttt{dag.rs:198}) lie en plus la chaîne
à une racine publiée. Deux re-prouveurs \emph{indépendants} existent : le CLI
\texttt{provenance-verify} (dépendances minimales) et \texttt{verify-mini} (\emph{zéro}
dépendance) — le sceptique peut re-vérifier avec un outil qui ne partage aucun code avec le
producteur.

\begin{cadre}
\textbf{Le refus vaut mieux qu'un faux vert.} Un bundle vide ne rend \emph{jamais} « VÉRIFIÉ » :
il est \textbf{REFUSÉ}. Un test de régression (\texttt{provenance-verify/src/main.rs:312--325})
épingle « bundle vide $\rightarrow$ exit 1, jamais 0 ». Et une erreur d'usage sort en \texttt{exit 2} — jamais un faux verdict
d'altération \texttt{exit 1} (\texttt{main.rs:116}). L'outil préfère dire « je ne peux pas
attester » que mentir par défaut.
\end{cadre}

\section{3 \;\textbar\; Couche 1 — l'artefact sans champ « conforme » (\texttt{reg-evidence})}
\texttt{derive()} (\texttt{reg-evidence/src/engine.rs:53}) \textbf{re-prouve le DAG d'abord} : si
la re-preuve échoue, il renvoie \texttt{RegError::NotProven} et « \emph{ne dérive JAMAIS
d'évidence sur du non-prouvé} » (\texttt{engine.rs:42}). Aucune évidence n'est construite sur des
records non re-prouvés.

La sortie est une \textbf{VUE}, pas un verdict. Le type \texttt{RegulatoryArtifact}
(\texttt{artifact.rs:83}) porte une \textbf{couverture} et les \texttt{node\_id} \emph{cités} — et
c'est tout :
\begin{lstlisting}
// artifact.rs — « outillé != conforme imposé par le TYPE » (l:2)
pub enum Coverage { Covered, Partial, OutOfDagScope }   // (l:49) couvert / partiel / hors-perimetre

pub struct RegulatoryArtifact {   // (l:83) « Aucun champ 'conforme' nulle part »
    // requirement, coverage: Coverage, cited_node_ids, semantics: note d'honnetete ...
    // -> il n'existe AUCUN champ ou ecrire « compliant » / « conforme »
}
\end{lstlisting}
Un développeur qui voudrait, « par pression commerciale, faire dire à l'artefact +système
conforme+, n'a structurellement nulle part où l'écrire » (\texttt{artifact.rs:8}) : le type
\emph{interdit} le glissement vers le verdict (\texttt{artifact.rs:57}). Deux gardes complètent
la couche : \texttt{PhantomCitation} (\texttt{engine.rs:37}) — une exigence ne peut pas citer un
\texttt{node\_id} absent du DAG re-prouvé (pas d'évidence fantôme) ; et surtout,
\texttt{MappingSet::new} \textbf{rejette un cadre qui ne déclare aucun trou}
(\texttt{mapping.rs:121}, garde \texttt{:130}) — « un cadre sans trou serait un \emph{vert
malhonnête} » (\texttt{mapping.rs:12}, test \texttt{mapping.rs:217}). Le template EU AI Act \emph{déclare} donc ses trous plutôt que de prétendre les
couvrir : art. 9 (gestion des risques), art. 10 (gouvernance des données), art. 11 + Annexe IV
(documentation technique), art. 17 (système qualité), art. 72 (surveillance post-marché) — tous
étiquetés \texttt{OutOfDagScope} (\texttt{templates.rs:73--81}). L'outil dit « je ne couvre pas
ça », au lieu de faire semblant.

\section{4 \;\textbar\; Les deux axiomes, gravés dans le type et testés au grep}
L'honnêteté n'est pas une promesse en prose — ce sont deux invariants \emph{exécutables} :

\smallskip
\textbf{Axiome 1 — outillé $\neq$ conforme.} Aucun champ de conformité dans l'artefact. Un test
(\texttt{artifact.rs:125--132}) sérialise l'artefact et \textbf{échoue le build} si le JSON
contient l'un des tokens \texttt{compliant}, \texttt{conforme}, \texttt{compliance} ou
\texttt{conformity}.

\smallskip
\textbf{Axiome 2 — intègre $\neq$ correct.} La ligne d'intégrité atteste « authentique et
non-altéré », \textbf{jamais} « correct » : elle dit atteste « \emph{ce qui s'est passé — pas la
justesse de la décision, ni une reproductibilité bit-à-bit} » (\texttt{audit-api/src/report.rs:74},
discipline \texttt{:6--7}). Un test (\texttt{report.rs:129}) interdit littéralement le token
« correct » dans le rendu.

\smallskip
Le point : on \emph{ne peut pas} sur-vendre par accident. Le type retire l'endroit où écrire le
verdict ; la crypto retire la possibilité de faire confiance au producteur ; les grep-tests
cassent le build si quelqu'un tente de faire passer « conforme » ou « correct ».

\section{5 \;\textbar\; Pourquoi c'est rare, et pourquoi un acheteur conformité y tient}
Refuser la coche verte est \emph{contre-intuitif} commercialement — c'est plus dur à vendre
qu'un « conforme ». Mais c'est la seule position honnête : la conformité est une conclusion
\emph{juridique} sur un déploiement précis, qu'aucun outil ne peut prononcer. La rareté n'est pas
d'y penser — c'est de rendre le refus \textbf{structurel} (pas de champ) plutôt qu'une note de bas
de page. Pour un acheteur conformité (achat, RGPD, AI Act), la proposition devient défendable :
\emph{« Certifié conforme ? Non, délibérément — reg-evidence n'a pas de champ conforme. Voici un
reçu d'évidence re-vérifiable ; votre juriste conclut. »} On ne lui remet pas un mensonge dont il
serait responsable — on lui remet une preuve qu'il peut re-vérifier lui-même, sans nous faire
confiance.

\section{6 \;\textbar\; Le périmètre honnête}
Ce que l'outil \textbf{n'affirme pas} : il ne prononce pas la conformité (outillé $\neq$
conforme), il n'atteste pas la justesse de la décision (intègre $\neq$ correct), il ne garantit
pas une reproductibilité bit-à-bit. Il atteste une seule chose : \emph{ce record est authentique
et non-altéré}, et voici l'évidence à laquelle chaque exigence se rattache — ou le trou déclaré
là où il n'y en a pas.

\smallskip
\textbf{Statut, sans fard.} La machinerie est \textbf{réelle et testée} : \texttt{provenance-dag},
\texttt{reg-evidence}, \texttt{audit-api}, \texttt{audit-console}, \texttt{provenance-verify},
\texttt{verify-mini} compilent et passent leurs tests, chaque garde ci-dessus étant épinglée par
un test nommé. La limite honnête d'aujourd'hui : le point d'entrée servi est un \textbf{binaire
local} dont le store est réhydraté depuis un fichier JSONL (ou provisionné vide) ; il n'ingère pas
encore du trafic de production capturé en direct (couche S1). Le \emph{moteur} est là et re-prouvable ;
le branchement bout-en-bout sur des données vivantes est l'étape suivante — et il ne changera pas
l'invariant : l'artefact n'aura toujours \emph{aucun champ} où écrire « conforme ».

\end{document}
